\clearpage\thispagestyle{empty}
\pdfbookmark[0]{Quellen zu den Beobachtungen der Victoria vom 31. Oktober 1850}{victoria740-sources}
\begin{center}\textbf{Quellen zu den Beobachtungen der Victoria vom 31. Oktober 1850}\end{center}
\noindent Diese KI-Übertragung ins heutige Deutsch folgt dem Auszug auf S.~478 der \href{https://gdz.sub.uni-goettingen.de/id/PPN236007467}{\textit{Werke}, Band~VI} (Göttingen, 1874). Carl Friedrich Gauß schrieb am 31. Oktober 1850 an Heinrich Christian Schumacher. Er teilte drei Beobachtungen der Victoria mit und beschrieb, was ihre Genauigkeit einschränkte.
\par\smallskip\noindent Der Brief erschien in den \textit{Astronomischen Nachrichten}, Band~XXXI, Nr.~740, \href{https://api.digitale-sammlungen.de/iiif/image/v2/bsb10049456_00157/full/1800,/0/default.jpg}{Sp.~305--306}. Die \href{https://api.digitale-sammlungen.de/iiif/image/v2/bsb10049456_00164/full/1800,/0/default.jpg}{Schlusszeile auf Sp.~320} nennt den 22. November 1850. Der Auszug der Werkausgabe endet nach den drei Beobachtungen. In der Zeitschrift folgen im Brief die von Goldschmidt berechneten Elemente und ein Vergleich mit Beobachtungen. KI-Übertragung und KI-Satz.
\clearpage
